\chapter{Herramientas CASE a utilizar}
\label{Herramientas CASE}

\section{Herramientas CASE en el contexto del proyecto} 

En  el desarrollo de software, la utilización de los diferentes tipos de herramientas CASE (Computer-Aided-Software-Engineering) resulta ser un factor muy importante para el éxito, organización y viabilidad de un proyecto tecnológico. Ya que estas herramientas no sirven solo para casos puntuales, son útiles en todo el ciclo de vida del desarrollo de software, desde el inicio con la planificación hasta el  final con la implementación.

Ahora bien, para el desarrollo de nuestro proyecto, el cual contempla una arquitectura compleja, con metodología ágil \textbf{Scrum}, las herramientas seleccionadas deben poseer algún carácter claramente colaborativo, flexible y orientado a la iteración continua. Dividiremos las herramientas en los 3 tipos estándar.

\section{Herramientas Upper CASE}

Estas son herramientas que se enfocan en las etapas iniciales del ciclo de vida del software, y por lo tanto su objetivo principal es ayudar en la recolección de requisitos, análisis del sistema y un diseño tanto arquitectónico como de interfaces antes de siquiera partir a desarrollar, entre las que encontramos especialmente útiles para nuestro proyecto se encuentran:

\begin{enumerate} 
	\item \textbf{Visual Paradigm}
	
	Esta aplicación especializada en hacer UML y BPMN, es particularmente útil en nuestro contexto, en donde la complejidad de los flujos de interacción en nuestra aplicación requiere de diagramas que validen la lógica del negocio antes de ser implementados, este entorno evalúa la sintaxis de los diagramas, evitando así errores conceptuales tempranos por los que no merece la pena perder tiempo en el futuro. 
	
	Se utilizará para generar el diagrama BPMN en el que se visualice como el cliente solicita un flete, el algoritmo de proximidad le notifica al transportista y como el administrador audita todo el proceso desde fuera. Asimismo la herramienta es elegida para generar los diagramas de casos de uso, diagramas de clases y diagramas de secuencia.
	
	\item \textbf{Figma}
	
	Esta herramienta es bastante conocida y no es por nada, es muy buena para diseñar el UI que posteriormente tendrá la aplicación y el UX, todo esto creando prototipos dinámicos interactivos, sirve para que el que posteriormente implementara no vaya a ciegas realizando la aplicación sin tener idea de cuál es su objetivo, ya que aquí puedes materializar los requerimientos funcionales que se hayan definido en pantallas navegables. Otra cosa muy buena que tiene es que es colaborativa y está basada en la nube, entonces no se requiere descargar nada para que los que estén encargados del diseño puedan trabajar mientras simultáneamente sus prototipos vayan siendo validados.
	
	Aquí se construirán todos los mockups que posteriormente serán detallados en la \textbf{Sección 10} del informe.
\end{enumerate}

\section{Herramientas Lower CASE}

Estas herramientas son las que toman el relevo una vez la arquitectura y diseño ya han sido definidos por las Upper CASE. Por su parte se utilizan en las fases de código, pruebas, optimización y despliegue de la aplicación, transformando así los modelos predefinidos en un producto tangible de software funcional. Además, en esta etapa se incluye la documentación de la aplicación.

\begin{enumerate}
	\item \textbf{Visual Studio Case}
	
	Plataforma fundamental que actúa como IDE ligero y extensible para la escritura del código. Es una herramienta ampliamente reconocida en la actualidad debido a la inmensa cantidad de extensiones que posee, lo que permite adaptar el entorno a las necesidades específicas de nuestro equipo para trabajar en prácticamente lo que sea. Para el desarrollo de las interfaces móviles y la plataforma web, nos facilita la vida al poder integrar diversas herramientas entre sí, permitiendo que los miembros del grupo trabajen de forma coordinada, mantengan un código ordenado y eviten errores de sintaxis al momento de programar la lógica del sistema de flete.
	
	\item \textbf{Postman}
	
	Recurso utilizado para la simulación y automatización de APIs, la naturaleza de nuestra aplicación es distribuida, la aplicación web y móvil deben comunicarse en tiempo real con el Backend. Esta comunicación se realiza a través de dichas APIs RESTful, las cuales deben ser probadas antes de que sean integradas en el Frontend, para así evitar fallos de rendimiento en un futuro.
	
	Para probarlas se pueden probar peticiones HTTP bajo diferentes escenarios, como el estrés de la carga de fotografías, simular la recepción continua de múltiples coordenadas GPS para actualizar el mapa en menos de 2,67 segundos.
	
	\item \textbf{TeXstudio} 
	
	TeXstudio es un IDE de LaTeX que proporciona un soporte moderno de escritura, como la corrección ortográfica interactiva, plegado de código y resaltado de sintaxis. Se usó principalmente para la escritura de este informe y para la documentación.
\end{enumerate}

\section{Herramientas I-CASE}

Finalmente las herramientas de tipo \textbf{I-CASE}, bastante importantes ya que están diseñadas como entornos globales que funcionan a lo largo de todo el ciclo de desarrollo. En un entorno como es Scrum, estas herramientas son cruciales para la gestión de Sprints, control de versiones y gestión de requerimientos.

\begin{enumerate}
	\item \textbf{Jira} 
	
	Instrumento utilizado para la gestión de proyectos ágiles, el método Scrum requiere de una administración continua de las tareas a realizar, para asegurar entregas de valor constantes. \textbf{Jira} proporciona el marco necesario para poder convertir las abstracciones del proyecto en elementos accionables, permitiendo medir el rendimiento del equipo para buscar comprender si se está realizando un buen trabajo y si hay algo mejorable. 
	
	Esta aplicación deberá ser gestionada por el Product Owner, quien la utilizaría para crear el Product Backlog, para poder generar posteriormente un Sprint Backlog.
	
	\item \textbf{Confluence} 
	
	Complementando la gestión de tareas, utilizaremos Confluence para la gestión de la documentación, está claro que un proyecto no puede sostenerse únicamente con código, a pesar de que en esta metodología la documentación no es tan extensa, sigue siendo importante, se requiere un lugar en donde todo el equipo pueda acceder para resolver dudas sobre las reglas del negocio del proyecto, aquí se detallarán las especificaciones más complejas, como fórmulas para calcular el porcentaje de una comisión, políticas para validar los documentos de los conductores, entre otras cosas. La mayor ventaja aquí es que se enlaza directamente con Jira.
	
	
	\item \textbf{GitHub}
	
	Esta plataforma es mundialmente conocida por ser el estándar en cuanto a control de versiones se trata. Trabajar con un grupo de siete u ocho estudiantes modificando los mismos archivos de código al mismo tiempo sería caótico sin un sistema que unifique el progreso. GitHub funciona como una herramienta de gestión colaborativa que centraliza nuestro trabajo y proporciona una trazabilidad completa de todos los cambios realizados. Utilizaremos estrategias de ramificación para que cada miembro pueda desarrollar una funcionalidad de forma aislada y luego solicitar una revisión por parte de un compañero. De esta forma nos aseguramos de que cualquier fragmento de código sea validado antes de integrarse a la rama principal, evitando que un error rompa el funcionamiento general de la aplicación de fletes.
	
\end{enumerate}
